\documentclass[12pt]{article}

\input{../../_assets/preambulo.tex}


\begin{document}

    % 1. Foto de fondo
    % 2. Título
    % 3. Encabezado Izquierdo
    % 4. Color de fondo
    % 5. Coord x del titulo
    % 6. Coord y del titulo
    % 7. Fecha

    
    \input{../../_assets/portada}
    \portadaExamen{ffccA4.jpg}{Ecuaciones\\Diferenciales II\\Examen XVIII}{Ecuaciones Diferenciales II. Examen XVIII}{MidnightBlue}{-8}{28}{2026}{}

    \begin{description}
        \item[Asignatura] Ecuaciones Diferenciales II.
        \item[Curso Académico] 2025/2026.
        \item[Grado] Doble Grado en Ingeniería Informática y Matemáticas.
        \item[Grupo] Único.
        \item[Profesor] Rafael Ortega Ríos.
        \item[Descripción] Convocatoria Ordinaria.
        \item[Fecha] 12 de junio de 2026.
        % \item[Duración] ---.
    
    \end{description}
    \newpage


    % ------------------------------------
    
    \begin{ejercicio}
        La solución maximal del problema
        \[ \dot{x} + \veps \cos x = 1, \quad x(0) = 0 \]
        se denota por $x(t; \veps)$.
        \begin{enumerate}[label=\alph*)]
            \item Prueba que la función $(t, \veps) \in \bb{R}^2 \mapsto x(t; \veps) \in \bb{R}$ está bien definida y pertenece a la clase $C^1(\bb{R}^2)$.
            \item Calcula $\displaystyle \del{x}{\veps}(2; 0)$.
        \end{enumerate}
    \end{ejercicio}

    \begin{ejercicio}
        Para cada $n \geq 2$ se considera la ecuación
        \[ x_n(t) = \pi + \frac{1}{n} x_n(0) + \int_0^t \operatorname{sen} x_n(s) \, ds, \quad t \in [-1, 1]. \]
        \begin{enumerate}[label=\alph*)]
            \item Demuestra que existe una única solución continua $x_n : [-1, 1] \to \bb{R}$.
            \item Calcula, si existe, $\lim\limits_{n \to \infty} x_n(t)$ para cada $t \in [-1, 1]$.
        \end{enumerate}
    \end{ejercicio}

    \begin{ejercicio}
        Dada una sucesión de funciones $f_n : [0, 1] \to \bb{R}$ de clase $C^1$, se sabe que $\{f_n\}$ converge uniformemente a $f$ y la sucesión de derivadas $\{f'_n\}$ converge uniformemente a $g$, ¿se puede asegurar que $g$ coincide con $f'$?
    \end{ejercicio}

    \begin{ejercicio}
        Se considera el problema de valores iniciales
        \[ \ddot{x} + 2\dot{x} + \sqrt{1 + \dot{x}^2} - 1 = 0, \quad x(0) = 0, \quad \dot{x}(0) = 1. \]
        Demuestra que existe una única solución maximal definida en toda la recta real.
    \end{ejercicio}

    \begin{ejercicio}
        ¿Es estable el equilibrio $x = 0$ para el sistema $x' = Ax$? Se supone
        \[ A = \begin{bmatrix} \nicefrac{-1}{2} & \nicefrac{-1}{2} & 3 \\ \nicefrac{1}{2} & \nicefrac{-1}{2} & 4 \\ 0 & 0 & -1 \end{bmatrix}. \]
    \end{ejercicio}

    \begin{ejercicio}
        Una sucesión de funciones $f_n : [0, 1] \to \bb{R}$ se dice \textit{equi-oscilante} si cumple $\forall \veps > 0 \ \exists \delta > 0 \ \exists N \in \bb{N} :$
        \[ |f_n(t) - f_n(s)| < \veps \quad \text{si} \quad |t - s| < \delta, \quad t, s \in [0, 1], \quad n \geq N. \]
        \begin{enumerate}[label=\alph*)]
            \item Prueba que si todas las funciones $f_n$ son continuas entonces una sucesión equi-oscilante es también equicontinua.
            \item Prueba que si $\{f_n\}$ es una sucesión equi-oscilante que converge uniformemente a $f$, entonces $f$ es continua.
        \end{enumerate}
    \end{ejercicio}

    \newpage
    \setcounter{ejercicio}{0} % Reiniciar contador de ejercicios
    % \noindent
    % \textbf{Solución.}

\end{document}